% SPDX-License-Identifier: CC-BY-4.0
\section{Introduction}\label{sec:introduction}

Polynomials that vanish to high order on a grid are a basic tool of the
polynomial method. Over the Boolean cube a natural extremal question asks how
small the degree of such a polynomial can be when it is required to vanish to
order \(k\) everywhere except at one point. Over \(\mathbb R\) this question was
answered by Sauermann and Wigderson~\cite{SW20}: the least degree of
\(P\in\mathbb R[x_1,\ldots,x_n]\) with \(P(\zero)\ne0\) and zeros of multiplicity
at least \(k\ge2\) at all other points of \(\{0,1\}^n\) is \(n+2k-3\). They observed
that the answer depends on the field and that characteristic two behaves
differently.

This paper determines the binary analogue. Throughout, multiplicity means Hasse
multiplicity of formal polynomials (Definition~\ref{def:mult}); in particular no
polynomial is reduced modulo \(x_i^2-x_i\), which would not preserve multiplicities.

\begin{definition}\label{def:admissible}
For integers \(n\ge1\) and \(k\ge1\), a polynomial
\(P\in\Ft[x_1,\ldots,x_n]\) is \emph{\(k\)-admissible} if
\(\mult_a(P)\ge k\) for every \(a\in\Ft^n\setminus\{\zero\}\) and
\(\mult_{\zero}(P)\le k-1\) \lean{claims/BinaryMultiplicityDegreeFormula.lean}.
Let \(D(n,k)\) be the least total degree of a \(k\)-admissible polynomial.
\end{definition}

The condition at the origin excludes \(P=0\) and allows any origin order below
\(k\), not only a nonzero value at the origin.
Bishnoi, Boyadzhiyska, Das and M\'esz\'aros~\cite[Question~4.3]{BBDM23} asked,
in the study of subspace coverings with multiplicities, whether
\(D(n,k)\ge2k-\lfloor k/2^{n-1}\rfloor\) when \(k\ge2^{n-2}\). This bound
follows for all \(n\) and \(k\) from the multiplicity Schwartz--Zippel lemma of Dvir,
Kopparty, Saraf and Sudan~\cite{DKSS13}; see Proposition~\ref{prop:cover-bound},
which was the content of an earlier note~\cite{Note}. That leaves open when the
bound is attained, and what \(D(n,k)\) is when it is not.

\subsection{Main result}

\begin{theorem}\label{thm:main}
For all integers \(n\ge1\) and \(k\ge1\) with \(k<3\cdot2^{n-1}\),
\[
 D(n,k)=2k+n-2-\lfloor\log_2k\rfloor
\]
\lean{claims/BinaryMultiplicityDegreeExtendedRange.lean}.
Moreover, for every \(n\ge1\) and \(k\ge1\) there is a \(k\)-admissible
polynomial of degree at most \(2k+n-2-\lfloor\log_2k\rfloor\)
\lean{claims/BinaryMultiplicityDegreeFormula.lean}.
\end{theorem}

The substance lies below \(2^{n-1}\), that is, for \(n\ge\lfloor\log_2k\rfloor+2\)
(Theorem~\ref{thm:main-formal}). For \(2^{n-1}\le k<3\cdot2^{n-1}\) the
Schwartz--Zippel bound already meets the construction (Corollary~\ref{cor:extended}).
Three features of the formula stand out.
\begin{itemize}
 \item \textbf{Where the Schwartz--Zippel bound is exact.} It equals the formula
 exactly when \(2^{n-2}\le k<3\cdot2^{n-1}\). For \(k<2^{n-2}\) it is never
 attained. For \(k\ge3\cdot2^{n-1}\) the two known bounds differ, and \(D(n,k)\)
 is open.
 \item \textbf{Each dimension costs one degree.} For fixed \(k\), \(D(n,k)-n\)
 is constant for \(n\ge\lfloor\log_2k\rfloor+2\). For example \(D(n,1)=n\),
 \(D(n,4)=n+4\), \(D(n,8)=n+11\) and \(D(n,15)=n+25\).
 \item \textbf{Comparison with the reals.} Sauermann and
 Wigderson~\cite{SW20} show that over \(\mathbb R\) the least degree with
 \(P(\zero)\ne0\) is \(n+2k-3\) for \(k\ge2\). Over \(\Ft\) our construction
 with \(P(\zero)\ne0\) (origin order \(\ell=0\)) has degree
 \(n+2k-2-s_2(k-1)\), which is \(s_2(k-1)-1\) lower. Allowing any origin order
 below \(k\) lowers the degree further, to \(D(n,k)\), which is
 \(\lfloor\log_2k\rfloor-1\) below \(n+2k-3\). Both savings are zero for
 \(k=2,3\), and the second requires \(n\ge\lfloor\log_2k\rfloor+2\).
\end{itemize}

We also prove the upper bound for each prescribed origin order in every
dimension (Theorem~\ref{thm:per-order}): for \(0\le\ell<k\) there is an
admissible polynomial with origin order exactly \(\ell\) and degree at most
\(n+2k-2-s_2(k-\ell-1)\), where \(s_2\) is the binary digit sum.

\subsection{Relation to previous work}

\begin{itemize}
 \item \textbf{Menezes~\cite{Men26}.} In a recent preprint, Menezes determines,
 over every field \(\F\), the least degree \(\delta_\F(n,k,\ell)\) of a polynomial
 with order at least \(k\) at the nonzero vertices of \(\{0,1\}^n\) and order
 exactly \(\ell\) at the origin, for \(k\ge2\) and \(n\ge k-1\). In characteristic
 two his formula is \(n+2k-2-s_2(k-\ell-1)\) \cite[Corollary~1.3]{Men26}.
 Minimizing over \(\ell\) gives Theorem~\ref{thm:main} when \(n\ge k-1\).
 Our contribution is the range \(\lfloor\log_2k\rfloor+2\le n<k-1\), which
 grows with \(k\). For \(k=2^j\) with \(j\ge3\), it is \(j+2\le n\le2^j-2\),
 while \(n\ge2^j-1\) is covered by~\cite{Men26}. The polynomials of our upper
 bound are Menezes' Catalan truncation specialized to \(\Ft\). His degree
 formula is stated for \(n\ge k-1\), the range of his theorem. We read his
 statements, not his proofs, so we do not claim that his argument is limited to
 that range. The telescoping identity of Lemma~\ref{lem:telescope} gives a
 short self-contained proof that the construction works in every dimension.
 Our lower bound is proved differently, by the one-dimension step of
 Theorem~\ref{thm:step} from the multiplicity Schwartz--Zippel lemma.
 \item \textbf{Sauermann and Wigderson~\cite{SW20}.} Their theorem is the real
 case described above. In characteristic two they give a polynomial of degree
 \(n+4\) for \(k=4\), below \(n+2k-3=n+5\), showing that the real formula
 fails there.
 \item \textbf{Dvir, Kopparty, Saraf and Sudan~\cite{DKSS13}.} Their
 multiplicity Schwartz--Zippel lemma is the only external input. It gives the
 base case \(n=\lfloor\log_2k\rfloor+2\) of the lower bound.
 \item \textbf{Alon and F\"uredi~\cite{AF93}.} For \(k=1\), Theorem~\ref{thm:main}
 says \(D(n,1)=n\) for \(n\ge2\). The lower bound is the case \(\F=\Ft\) of
 their Theorem~5, which our argument reproves.
\end{itemize}

Two web searches, for high-order vanishing on the hypercube in positive
characteristic and for punctured multiplicity Nullstellens\"atze over
\(\Ft\), found no statement for the range \(n<k-1\) other than the results
above. A bounded search does not establish that none exists.

\subsection{Proof outline}

\begin{itemize}
 \item \textbf{Lower bound (Section~\ref{sec:lower}).} Take an admissible
 \(P\) in \(n\) variables. After an invertible linear change of coordinates
 in which the lowest homogeneous part at the origin is not divisible by
 \(x_n\), the polynomial \(S=P(x',0)+P(x',1)\) is admissible in \(n-1\)
 variables, has the same origin order, and has degree at least one lower
 (Theorem~\ref{thm:step}). Iterating down to \(m=\lfloor\log_2k\rfloor+2\),
 where the multiplicity Schwartz--Zippel lemma already gives degree \(2k\),
 yields \(D(n,k)\ge2k+n-m\).
 \item \textbf{Upper bound (Section~\ref{sec:upper}).} With
 \(y=x^2+x\) we have \((1+x)+\sum_{j<K}y^{2^j}=(1+x)^{2^K}\) over \(\Ft\).
 Expanding \(\prod_i(1+x_i)^{2^K}\) this way and discarding the terms of
 cost at least \(s\) leaves a polynomial \(g_s\). It still vanishes to order
 \(s\) off the origin, takes the value \(1\) at the origin, and has small
 degree (Theorem~\ref{thm:truncated-product}). Multiplying by
 \(y_1^\ell\) sets the origin order (Theorem~\ref{thm:per-order}).
\end{itemize}

\subsection{Formal verification}

Theorem~\ref{thm:main}, every lemma it uses, and the multiplicity
Schwartz--Zippel lemma itself are proved in Lean~4 with Mathlib. The proofs
use only Lean's standard axioms. Each statement in this paper carries a
[Lean] link to its formal proof at a fixed repository commit.
Appendix~\ref{sec:verification} maps statements to declarations and gives the
commands to check them. The external input is stated without proof in the
text, with a citation, but its formal proof is linked like every other
statement, so the formal verification has no unchecked premise. Formal
verification covers the statements under the linked definitions. It does not
cover the prose, the comparison with the literature, or the fidelity of the
definitions to the informal notions, although the definitions are short
(Definitions~\ref{def:mult} and~\ref{def:admissible}).
